\section{Finite algebraic certificates and real descent}\label{app:foundations}
This appendix supplies the full family argument used in
Lemma~\ref{lem:spread} and Theorem~\ref{lem:relative}. A variety in the
algebraic construction is a reduced separated scheme of finite type over
a characteristic-zero field. Smoothness is geometric smoothness. For
an effective input the ground constants belong to a specified finite real
algebraic extension of $\Q$; generic fibres are presented over its
finitely generated function fields. No algorithmic assertion about an
arbitrary computable field is intended. Real realizations are subsequently
taken in $\R$, with the original sign and branch conditions retained.

An algebraic coordinate presentation consists of finitely many affine
coordinate rings, localized by specified nonzero elements, transition
morphisms and their inverse identities, and a coordinate map whose
relative Jacobian is invertible. An adapted presentation also lists each
divisor equation, its expression as a coordinate times a unit, and the
monomial identities for the functions. Projective maps are recorded by
homogeneous equations or by their finite blow-up sequences. These are
finite polynomial data, not a numerical collection of coordinate boxes.

\subsection{The Rees algebra and a finite diagram}
\begin{lemma}[One localization for all powers of an ideal]\label{lem:rees}
Let $A$ be a noetherian domain, $B$ an algebra of finite type over $A$,
and $I\subset B$ a finitely generated ideal. There is $a\in A\setminus\{0\}$
such that, for every $A_a$-algebra $A'$, with $B'=B_a\otimes_{A_a}A'$,
\[
 I_a^n\otimes_{A_a}A'\ \xrightarrow{\ \sim\ }\ (I_aB')^n
 \quad(n\ge0),\qquad
 \left(\bigoplus_{n\ge0}I_a^n\right)\otimes_{A_a}A'
      \ \cong\ \bigoplus_{n\ge0}(I_aB')^n.
\]
The first isomorphism identifies both sides as submodules of $B'$.
Consequently the blow-up in $I_a$ commutes with this base change. One
localization suffices for any specified finite sequence of centre ideals
on finitely many affine charts of successive blow-ups.
\end{lemma}
\begin{proof}
The associated graded algebra
$\operatorname{gr}_I B=\bigoplus_{n\ge0}I^n/I^{n+1}$ is of finite type
over $A$: it is generated by $B/I$ in degree zero and the images of a
finite set of generators of $I$ in degree one. Since $A$ is noetherian,
these algebras are finitely presented. Apply generic freeness
\cite[Tag 051S]{Stacks} to $B$ as a module over itself and to
$\operatorname{gr}_I B$ as a module over itself. After inverting the
product of the two resulting elements, both are flat over $A_a$.
Each graded piece is a direct summand as an $A_a$-module, hence flat.
In particular $B_a/I_a$ is flat. Induction using
\[
 0\longrightarrow I_a^{n-1}/I_a^n\longrightarrow B_a/I_a^n
       \longrightarrow B_a/I_a^{n-1}\longrightarrow0
\]
shows that $B_a/I_a^n$ is flat for every $n\ge1$.
Tensoring $0\to I_a^n\to B_a\to B_a/I_a^n\to0$ with $A'$ is
therefore exact on the left. Its image in $B'$ is the ideal generated
by products of $n$ images of elements of $I_a$, namely $(I_aB')^n$.
The isomorphisms respect multiplication and grading, proving the Rees
algebra assertion. Relative Proj commutes with base change; this proves
the blow-up assertion. For a finite sequence, apply the argument to the
actual centre ideal on every affine chart at each stage and multiply the
finitely many localizing elements. Previously proved identities survive
further localization. No countable intersection of open sets is used.
\end{proof}

\begin{lemma}[Spreading the entire stratifying diagram]\label{lem:finite-diagram}
Let $A$ be a finitely generated integral algebra over a characteristic-zero
field $k$, and let $K=\operatorname{Frac}(A)$. Suppose $Y\to\operatorname{Spec}A$
is of finite type and a generic certificate consists of a reduced closed
chain $Y_{0,K}=Y_{K,\mathrm{red}}\supset\cdots\supset Y_{q,K}=\varnothing$,
projective $Z_{j,K}\to Y_{j,K}$, and specified inverse isomorphisms above
$E_{j,K}=Y_{j,K}\setminus Y_{j+1,K}$. Then, after replacing $A$ by $A_a$
for one nonzero $a$, the chain can be chosen with $Y_0=Y_{\mathrm{red}}$,
all its members flat with geometrically reduced fibres, and the stated
isomorphisms hold above the actual complements
$E_j=Y_j\setminus Y_{j+1}$. In particular they hold above $(E_j)_s$ on
every fibre. The union of these complements is the whole reduced fibre;
no constancy of the number of irreducible components is required.
\end{lemma}
\begin{proof}
First replace $Y$ by its reduction. Take reduced closures of the generic
$Y_{j,K}$ inside $Y$ and use $Y_q=\varnothing$. Clear denominators in
the projective equations and affine presentations of the generic maps.
All schemes involved are of finite presentation over the noetherian
base. The finite-presentation limit lemma
\cite[Tag 01ZM, Lemma 32.10.1(1)--(3)]{Stacks} descends objects,
morphisms, and equalities of morphisms. Apply it to the inverse maps on
the generic complements, their inclusions, and both inverse identities.
The complements used on the model are $Y_j\setminus Y_{j+1}$; their
generic fibres are the prescribed complements. The descended inverse
maps and equalities therefore give isomorphisms on these actual open
subschemes after a further localization, not merely on unspecified
subsets with the same generic points. Finite affine covers and their
intersection diagrams are included in this application of the limit
lemma. Thus chart coverage, as well as the formula for each map, descends.

Each reduced generic member is geometrically reduced because $K$ has
characteristic zero. Generic freeness makes each member flat on a dense
open set. The generic-fibre reducedness theorem
\cite[Section 37.26, Lemma 37.26.4]{Stacks} then permits a further common
shrink so that all fibres of these finitely many members are geometrically
reduced. This is an imposed property, not an optional replacement of
scheme structure after specializing. Finally, on every fibre the sets
$Y_{j,s}\setminus Y_{j+1,s}$ form a disjoint partition of $Y_{0,s}$
because the chain ends in the empty set. For any point there is a unique
index at which it leaves this finite descending chain. The corresponding
isomorphism covers that point, including points on a special fibre
component. This argument neither identifies fibre components with generic
components nor assumes their number to be constant.
\end{proof}

\subsection{Relative smoothness, normal crossings, and exceptional sets}
We now prove Lemma~\ref{lem:spread} with a precise list of the further
conditions. Separate the finitely many open-and-closed components of each
smooth generic source according to the identically zero functions. These
components and their idempotent identities are part of the finite diagram.
On a component on which a function is declared zero, descend that zero
identity. On every other component principalize the product of the
remaining functions, recording all factors and units.

The absolute input is characteristic-zero resolution and principalization
\cite{BM97,BM08}. More precisely, the variety resolution of
\cite[Theorem 1.1]{BM08} supplies the smooth sources; the marked-ideal
resolution of \cite[Theorem 1.3]{BM08}, applied with mark one, supplies
principalization together with the normal-crossing exceptional boundary.
A product which is a monomial times a unit makes each of its factors a
monomial times a unit on the regular local rings. This use is over the
field $K$. Compatibility with arbitrary special fibres is supplied by
Lemma~\ref{lem:rees}, not by a claim that functoriality for smooth maps
implies arbitrary base-change compatibility.

Spread the finite sequence of centres and its adapted charts. The
following finite conditions are imposed before retaining the open base.
The source and centre smooth loci are the loci on which their finite
presentations satisfy the relative smoothness criterion. On a smooth
source the intersections of the designated divisor equations must be
smooth of their prescribed relative codimensions; in adapted coordinates
this says that the appropriate relative differentials are independent.
This condition is imposed for every subset of the finite divisor list.
Each coordinate chart must be etale, each declared inverse unit must
multiply its unit to one, and the finite charts must cover their source.
All scheme-theoretic identities in these statements are spread by the
finite-diagram lemma. Their remaining rank conditions are open algebraic
conditions; complements of chart unions are constructible algebraic sets.

For clarity, the discarded algebraic parameter set is the union of the
closures of images of the following bad sets, together with the
localizing hypersurfaces already used:
\begin{enumerate}
\item the original irreducible components not dominating the current
reduced integral base;
\item nonsmooth points of a source or centre, and failures of the
specified relative dimension or centre codimension;
\item failures of smoothness or codimension of any designated divisor
intersection, or failure of an adapted etale Jacobian;
\item the complement of the chosen finite source charts, and any zero
of a function declared invertible;
\item failures of a descended zero identity, a chart transition, or an
inverse isomorphism on an actual complement in the closed chain.
\end{enumerate}
For items (4)--(5) the bad set is empty once the relevant identities and
coverage have descended; it is harmless to include it explicitly in the
certificate. Every bad set in this list has empty generic fibre: this is
respectively dominance, generic smoothness in the given certificate,
generic normal crossings, generic coverage and invertibility, or the
specified generic identities. Images are constructible by Chevalley's
theorem \cite[Tag 054K]{Stacks}. A constructible subset of an integral
noetherian space which is dense contains a nonempty open subset and hence
the generic point. Therefore the closure of each of these images is
proper. This justification applies also when $Y\to S$ is not proper.

After their removal, apply Lemma~\ref{lem:rees} on every spread blow-up
chart. The resulting smooth sources and centres, relative normal crossings,
unit identities and finite chart covers hold on every retained fibre.
A nonzero monomial cannot specialize to the zero function on a fibre
component: its unit is invertible and its finitely many divisors each have
relative codimension one. Their union cannot contain a component of the
smooth fibre. An identically zero generic function remains zero by its
descended identity. The open-and-closed source decomposition prevents a
change of these flags through an unlabelled component intersection.
Flatness of the closed chain and smoothness of the sources exclude
additional vertical components on the retained open set. Coverage of any
special fibre components is, more directly, the partition argument in
Lemma~\ref{lem:finite-diagram}. This proves all fibrewise assertions of
Lemma~\ref{lem:spread}.

Each discarded algebraic subset has strictly smaller dimension than its
current integral base. Replace it by its reduced irreducible components
and reconstruct the entire generic certificate on each component. There
are finitely many components at each step and depth at most $\dim S+1$.
This gives finite termination. The exceptional list is not required to be
minimal. Importantly, a change of a \emph{real} accessibility flag need
not lie in a proper algebraic subset: two full-dimensional real chambers
can have different flags. Such changes are treated by finite real
quantifier-elimination partitions below, not by this dimension argument.

\subsection{Exact real algebraization}
\begin{lemma}[The square-slack cover with all branch conditions]\label{lem:slack-exact}
Let $S$ be a semialgebraic base and $\Gamma\to S$ a semialgebraic proper
map realized as a graph in $S\times\R^N$. The graph includes all auxiliary
functions, sign choices and denominator exclusions. After a finite base
refinement, it admits finitely many algebraic square-slack lifts whose
real points over $S$ project exactly onto $\Gamma$. Projection is proper
over compact subsets of $S$. Every later map factoring through these
lifts projects only to points of $\Gamma$.
\end{lemma}
\begin{proof}
Let $C$ be the ordinary semialgebraic closure of $\Gamma$ in the ambient
real algebraic space. Properness implies that $\Gamma$ is relatively
closed in $S\times\R^N$: a convergent sequence over a convergent base
sequence lies over a compact subset of $S$ and has a graph limit there.
It follows that $C\cap(S\times\R^N)=\Gamma$. A closed semialgebraic
set is a finite union of basic closed semialgebraic sets
\cite[Theorem 2.7.2]{BCR98}. For a piece
\[
 C_l=\{x:f_{l,a}(x)=0,\ p_{l,b}(x)\ge0\}
\]
use the algebraic set
\[
 Y_l=\{(x,s): f_{l,a}(x)=0,\ p_{l,b}(x)=s_b^2\}.
\]
Every real point of $Y_l$ projects into $C_l$, and each point of $C_l$
has a real lift. Over $S$, the union therefore projects exactly to
$\Gamma$. Over a compact base subset the graph is compact. The slack
coordinates are bounded by the square roots of maxima of the finitely
many polynomials there, and the lift is closed; hence the lifted real
set is compact. This proves properness and the last assertion.

Strict exclusions are encoded in the complete graph formula before this
construction. For a denominator excluded everywhere on that compact graph,
$d\ne0$ may also be retained by adjoining $w$ with $dw=1$, together
with any required sign of $d$. Nonnegative square-root branches are
specified in the graph, not selected after taking an algebraic closure.
If adjoining such coordinates destroys properness on the initial base,
first apply Hardt triviality to the enlarged graph, whose individual
fibres are compact, and refine the base. A trivialization with compact
fibre makes the restricted graph map proper. Closure is taken only
\emph{after} this step. Consequently no point with a vanishing excluded
denominator is reintroduced over the retained stratum. General Boolean
branch conditions are instead kept in the complete graph formula: one
does not split a compact fibre into nonclosed pieces, adjoin unbounded
reciprocals on those pieces, and then assume their fibres are compact.
The equality $C\cap(S\times\R^N)=\Gamma$ already preserves all such
conditions without that operation.
\end{proof}

\begin{lemma}[Real branch selection and simultaneous descent]\label{lem:real-descent}
The root labels and Nash graph functions used in the cluster construction
can be carried on finitely many selected real Nash branches of algebraic
base covers. On a finite semialgebraic refinement, the proper resolved
square-slack sources cover exactly the original graph in every fibre.
Descent by root permutations preserves the entire joint coefficient and
root relation, not merely its scalar order list.
\end{lemma}
\begin{proof}
Refine a Nash stratum by the regular locus of each integral component of
its algebraic closure. Its lower-dimensional singular parts are treated
recursively. A Nash function is algebraic over the coordinate function
field. Adjoin the finitely many graph coordinates needed here to that
field. The extension is finite and separable in characteristic zero.
Choose equations and invert their leading coefficients and discriminants;
a further localization gives a finite etale algebraic cover. Include the
original graph formula to select the intended real branches. Finite
semialgebraic cell decomposition makes each selected graph a Nash section
and separates its alternatives. Complex sheets and real sheets failing
the original formula are not retained. Thus monodromy is handled by base
refinement, rather than by globally labelling a nontrivial cover.
Exceptional leading coefficients and discriminants are handled by the
reduced lower-dimensional base recursion. A finite union of selected
sections has compact inverse image over compact subsets of a stratum.

Apply Lemma~\ref{lem:slack-exact} on each selected branch. For a real
point $y$ of a lifted fibre, choose the index with
$y\in Y_{j,s}\setminus Y_{j+1,s}$ in the spread singular-locus chain.
The inverse isomorphism of Lemma~\ref{lem:finite-diagram} supplies a real
point of its smooth source above $y$. A blow-up of a smooth real variety
in a smooth centre of positive codimension $c$ has fibre
$\R\mathbb P^{c-1}$ at a real centre point and is an isomorphism away
from the centre. In local smooth coordinates the centre is a coordinate
subspace, which proves both assertions directly. Such a blow-up therefore
does not lose real points. Composing the finitely many principalizations
and taking the finite union over $j$ proves real surjectivity. A real
singular point with no top-dimensional real branch is covered by a later
member of the chain. Projectivity of these maps and properness of the
slack projection prove properness of the real cover.

Conversely every real resolved point factors through a real slack point
on the selected branch, and hence projects to the original graph by
Lemma~\ref{lem:slack-exact}. This proves equality, not just inclusion,
of the covered real fibres. A root-label permutation simultaneously
permutes cluster indices, their coefficient coordinates and root arrays.
The equations defining $\Psi$ and the fixed-centre formula
\eqref{eq:closure-formula} are equivariant for this simultaneous action.
A finite union commutes with taking a fibrewise closure; projection to the
finite quotient therefore preserves the joint relation and its root
image. Descent never pairs one coefficient from one sheet with another
coefficient from a different sheet. The bottleneck metric is invariant
under these permutations.
\end{proof}

\subsection{Real accessibility and completeness of the order list}
\begin{lemma}[Real-accessible divisors]\label{lem:real-accessibility}
On a resolved square-slack source with the above relative adapted
coordinates, accessibility of a divisor is a first-order real condition
on the parameter. Every recorded divisor has an admissible real
transversal attaining its monomial orders. Every real degeneration
relevant to a finite coefficient order is controlled by one of the
recorded divisors. Thus the minimum in Corollary~\ref{cor:orders} is both
an upper-inequality exponent and an attained exponent.
\end{lemma}
\begin{proof}
In a chart with source equations $E(\theta,y)=0$, chart exclusions
$D(\theta,y)\ne0$, and adapted divisor coordinates $y_1,\ldots,y_l$,
the flag for $i$ with $a_i>0$ is the formula
\begin{equation}\label{eq:accessibility-flag}
 \exists y:\ E(\theta,y)=0,\quad D(\theta,y)\ne0,\quad
 y_i=0,\quad\prod_{h\ne i,\,h\le l}y_h\ne0,
 \quad \det J_{\rm et}(\theta,y)\ne0.
\end{equation}
Inverse units and the selected real base formula are included in $E,D$.
Charts are the ones meeting the inverse image of $G=0$. Since the full
real slack fibre already represents the admissible graph, there is no
additional unrecorded sign restriction on source directions. At a point
satisfying \eqref{eq:accessibility-flag}, the real inverse function
theorem gives a transversal with $y_i=t$ and other adapted coordinates
fixed. For small nonzero $t$, $G>0$, and every point projects into the
original graph. The unit bounds give $G\asymp |t|^{a_i}$ and
$H\asymp |t|^{b_i}$ for each nonzero $H$ on that source component.

A real degeneration with $G\to0$ has a lifted accumulation point above
$G=0$, by properness. Some adapted chart contains that point. If it lies
on several divisors, relative normal crossings allows a small real
perturbation \emph{on any one of them} avoiding all the others. Hence
each divisor of positive $G$ order meeting this real chart has a point
of the form \eqref{eq:accessibility-flag}. In a smaller box, with
$\beta$ the minimum of its recorded ratios, $b_i\ge\beta a_i$ for
$a_i>0$ and $b_i\ge0$ otherwise. The monomial identities and positive
unit bounds imply $H\le C G^\beta$ throughout the box, including its
intersections. Equivalently, along a lifted arc the order ratio is a
weighted average of the $b_i/a_i$ with nonnegative additional terms
from indices with $a_i=0$; it cannot be smaller than their minimum.
A finite cover near the compact exact fibre gives this inequality for
all competitors. The transversal for a minimizing divisor attains it.
If $G$ is identically zero on a source with real points, the exact-target
condition gives $H=0$ there. A smooth fibre component with a real point
has a Zariski-dense real open subset. Thus a nonzero monomial $H$ cannot
occur on such an exact real component; real-empty components contribute
no order. This handles zero flags as well as the nonzero ratios.

Finally \eqref{eq:accessibility-flag} is a finite polynomial formula.
Quantifier elimination partitions the base by all its truth values.
These semialgebraic pieces may have full dimension. Nash refinement
makes the choices compatible with the selected base branches. This
partition is finite without an algebraic dimension-drop assertion.
\end{proof}

\begin{proof}[Completion of Theorem~\ref{lem:relative}]
Stratify the enlarged graph and the exact cluster data as in
Lemma~\ref{lem:slack-exact}, and carry labels as in
Lemma~\ref{lem:real-descent}. Over an integral algebraic base resolve the
reduced generic fibre, then its reduced singular locus, and continue until
the empty set. Dimension decreases in this chain. Principalize on its
smooth sources the product of the nonzero functions. Lemmas
\ref{lem:rees} and \ref{lem:finite-diagram}, followed by the explicit
bad-locus construction, spread this certificate over a nonempty open
base with all fibrewise properties. Repeat on the finitely many reduced
exceptional components. Lemma~\ref{lem:real-descent} proves exact real
coverage after descent, and Lemma~\ref{lem:real-accessibility} proves
the accessibility assertions and the completeness of the orders.

For a fixed fibre, cover the compact inverse image of $G=0$ by finitely
many smaller real coordinate boxes with positive lower and finite upper
unit bounds. Properness implies that a sufficiently small observation
neighbourhood is covered: otherwise a sequence with $G\to0$ outside
the chosen source neighbourhood would have a lifted accumulation point
above $G=0$ outside it. The same argument works above a compact base
subset, since the selected label graph, the slack lift and the projective
resolved maps are proper there. Alternatively their positive radii and
unit bounds can be selected semialgebraically and made continuous by
finite refinement. Their extrema on a compact parameter subset give
common neighbourhoods and bounds. Fibrewise proper source maps must not
be confused with the smaller, nonproper coordinate boxes. This proves
all existence and compact-uniform assertions. The effective construction
for algebraic input is detailed at the end of Appendix~\ref{app:effective}.
\end{proof}

\section{Finite support orders, uniform decay, and effective output}\label{app:effective}
The metric and effective parts can be proved without making a
parameter-dependent choice of a Puiseux branch. Only the ordinary
one-variable algebraic-branch theorem is used to prove that a finite
candidate list is exhaustive. Selection from that list is then a
first-order decision problem.

\begin{lemma}[A finite support test for the order]\label{lem:support-order}
Let $e(\theta,t)\ge0$ be a semialgebraic function, defined for
$0<t<\tau(\theta)$, where $\tau>0$ is semialgebraic, and suppose
$e(\theta,t)\to0$ for fixed $\theta$. A quantifier-free graph formula
supplies a finite computable set $\mathcal B\subset\Q_{>0}$ such that
for every $\theta$ either $e(\theta,t)$ is eventually zero or exactly
one $b\in\mathcal B$ satisfies
\begin{equation}\label{eq:order-decision}
 \exists c,C,a>0\ \forall t:\quad
 0<t<a\ \Longrightarrow\ c t^b\le e(\theta,t)\le C t^b,
 \qquad a<\tau(\theta).
\end{equation}
Each of these alternatives is an effective semialgebraic condition for
algebraic input. No real exponent is quantified over.
\end{lemma}
\begin{proof}
Eliminate quantifiers in the graph formula and let $P_1,\ldots,P_N$
be its finite list of polynomials in $(\theta,t,y)$. Fix $\theta$ and
discard those specializing to the zero polynomial in $(t,y)$. At a graph
point some remaining polynomial must vanish. Otherwise all nonzero
polynomial signs are unchanged on a small two-dimensional neighbourhood;
the zero specialized polynomials impose constant truth values there.
The graph formula would then contain that neighbourhood, contradicting
single-valuedness. For each remaining $P_i$, the set
$\{t:P_i(\theta,t,e(\theta,t))=0\}$ is semialgebraic in one variable.
These finitely many sets cover a punctured interval. After their finitely
many endpoints are removed, at least one vanishing equation holds on an
entire interval $(0,a)$. A polynomial independent of $y$ cannot vanish
there unless it was one of the discarded zero polynomials.

By one-variable semialgebraic decomposition, on a smaller interval the
function is either identically zero or positive, continuous, and a single
real algebraic branch of such an equation. The ordinary algebraic
Puiseux theorem gives, in the positive case,
$e(\theta,t)=a_\theta t^{b_\theta}+o(t^{b_\theta})$ with
$a_\theta>0$ and rational $b_\theta>0$; see \cite{BCR98,BPR06}.
Write its vanishing equation as
$\sum p_{ij}(\theta)t^i y^j=0$. Inserting the first term shows that
$\min(i+b_\theta j)$ is achieved at least twice among nonzero
specialized coefficients. Otherwise its unique lowest term could not
cancel. Therefore
\[
 b_\theta=(i-i')/(j'-j)>0
\]
for two distinct support points of one of the original $P_i$.
Take $\mathcal B$ to be all such positive ratios. Specialization may
remove support points but cannot add a pair to this list. This proves
finiteness uniformly in $\theta$, without following algebraic branches
through parameter specializations.

Eventual zero is the formula
$\exists a>0\ \forall t\,(0<t<a\Rightarrow e(\theta,t)=0)$, with
$a<\tau(\theta)$. For each candidate, \eqref{eq:order-decision} is a
first-order polynomial formula: replace $t^{p/q}$ by its unique positive
root, or put $t=s^q$ throughout. Quantifier elimination decides it.
The actual branch expansion supplies one true candidate; two distinct
powers cannot be bounded above and below by positive multiples of each
other as $t\downarrow0$. Exactly one candidate is therefore true in
the nonzero case. This also tracks the branch selected by the original
Boolean graph formula, including all specialized coefficients and
inequality faces, rather than selecting a root of a polynomial without
checking graph membership.
\end{proof}

\begin{proof}[Full proof of Proposition~\ref{prop:decay}]
Use Lemma~\ref{lem:support-order}. If the candidate set is nonempty,
choose a positive rational $\nu<\min\mathcal B$; if it is empty use
$\nu=1$. On every parameter fibre $e(\theta,t)\le t^\nu$ eventually,
including the eventual-zero case. The set of pairs $(\theta,a)$ with
$0<a\le\min(1,\tau(\theta))$ and this inequality for all $0<t<a$
is semialgebraic and has a nonempty positive fibre at every parameter.
Semialgebraic choice selects a positive $a=\tau_0(\theta)$, not a
possibly vanishing limit value at the boundary of a stratum. A finite
Nash stratification makes this selected function continuous. Its minimum
on a compact subset of a stratum is strictly positive; the estimate is
uniform there. The exponent may be smaller than the sharp order, which
is exactly what is needed for the uniform remainder.
\end{proof}

\begin{lemma}[Nash refinement of graph functions]\label{lem:nash-graph}
Finitely many semialgebraic functions on a semialgebraic set admit a
common finite Nash stratification on which each function is Nash.
For algebraic presentations this refinement is effective.
\end{lemma}
\begin{proof}
Take their combined graph. Semialgebraic Nash stratification, refined by
the ranks of the coordinate projection and by a compatible stratification
of its image, is finite; these are the usual graph and rank-stratification
constructions in \cite{BCR98,BPR06}. On a graph stratum the projection is
injective. If its constant rank were less than the dimension of that
stratum, the constant-rank theorem would give a positive-dimensional
local fibre, contradicting injectivity. The projection thus has full
rank onto its image stratum. Its local inverse is analytic by the inverse
function theorem and is semialgebraic because its graph is semialgebraic;
hence it is Nash. Finitely refine images to separate overlapping strata
and repeat on their lower-dimensional boundaries. This proves the claimed
finite partition and gives the precise reason that graph stratification,
rather than topological triviality alone, yields Nash functions.
Effective semialgebraic decomposition and rank tests give the algebraic
version of the construction.
\end{proof}

The application to Theorem~\ref{thm:uniform} is now direct. The polynomial
matching formula \eqref{eq:matching-formula}, with two quantified
inclusions and minimality over the distance, defines the Hausdorff error
as a semialgebraic function of $(\theta,t)$. Theorem~\ref{thm:joint}
proves its pointwise convergence to zero. Proposition~\ref{prop:decay}
gives a compact-uniform power estimate. Lemma~\ref{lem:nash-graph}
applied to the diameter and the positive separation and radius functions
gives the asserted Nash refinement. Hardt triviality is used separately
for the topological type of the leading fibres. The inequality
$|\diam A-\diam B|\le2d_H(A,B)$ proves the diameter remainder.

\subsection{An explicit terminating spectral algorithm}
Here is a proof of Theorem~\ref{thm:effective} with its data representations
and branch-selection step made explicit. The input is the polynomial
formula of Definition~\ref{def:input}; each algebraic constant has its
integer polynomial and rational isolating interval. Arithmetic is in a
primitive-element representation of their finite real algebraic field.
All quantified formulas below are subjected to real quantifier elimination
as in \cite{BPR06}. Finite root permutations are literal finite
disjunctions, and complex coordinates are pairs of real coordinates.

First decide exact identification by the formula that all $x,x'$ with
$F(x)=F(x')=P$ have equal target coefficients. On this locus construct
the cluster graph using its polynomial factorization equations and the
specified isolating conditions. For each $H=|c|^2$, form the graph of
\[
 m_H(s)=\max\{H(x):x\in\Gamma,\ G(x)\le s\}
\]
by the conjunction of the universal upper-bound formula and an existential
attainment formula. Compactness gives the maximum; compactness and
$\{G=0\}\subset\{H=0\}$ give $m_H(s)\to0$. Apply quantifier
elimination to this graph. Construct its finite support-pair list as in
Lemma~\ref{lem:support-order}, test eventual zero, and test
\eqref{eq:order-decision} for each rational candidate. The unique true
candidate is $\beta_H$. This is a finite algorithm for selecting the
\emph{actual maximum-function branch}. It needs no unspecified command
to isolate an algebraic branch at the boundary. If all maxima are eventually
zero the target is locally rigid; otherwise set
$\alpha=\min_H\beta_H/k_H$ over the nonzero maxima.

For a family, keep $\theta$ free in exactly these tests. Their finitely
many truth sets give the parameter pieces and include all zero
specializations. Nash-refine them using Lemma~\ref{lem:nash-graph}.
The supports are finite before specialization, so this procedure still
terminates without a bound on a real-valued exponent.

Next substitute $\alpha=p/q$ in the fixed-centre formula
\eqref{eq:closure-formula}. Elimination computes a quantifier-free
presentation of $\cI_0$ and its projection $\cC_0$. Adjoin the labelled
roots by elementary symmetric identities to obtain $\mathcal L_0$.
Use \eqref{eq:matching-formula} to express that all pairs of elements
have distance at most $C$ and that some pair attains $C$. Compactness
makes this a formula for one positive real number. Univariate real-root
isolation applied to its eliminated formula selects that number over the
input field. Taking a field norm gives a nonzero integer polynomial
vanishing at $C$; renewed root isolation and the same formula select a
rational isolating interval for the intended root. The same graph formulas
with free parameters compute the family of constants, and the preceding
Nash and decay constructions compute a compatible refinement and an
admissible rational remainder exponent. No complexity bound is claimed.

\subsection{Effective geometric certificates without an implicit oracle}
The geometric output can also be given a terminating, though generally
inefficient, specification. Over the field just described, enumerate
finite algebraic diagrams and finite semialgebraic base partitions in
increasing finite encoding length, refining a fixed semialgebraic
triviality partition of the enlarged graph with compact fibres. Properness
on those base pieces follows from that trivialization and is preserved
under further base restriction. A candidate consists of square-slack
presentations, reduced closed chains with inverse-map identities, smooth
sources with specified projective maps, finite centre ideals and adapted
etale charts, integer monomial orders, inverse units and selected real
branch formulas. Verify polynomial identities and centre blow-up charts
by ideal arithmetic; verify relative smoothness, chart ranks and divisor
intersections by their Jacobian criteria. Properness is certified by the
specified projective construction and the square-slack argument, not by
numerical bounds on open charts. Verify partition coverage, fibrewise
real surjectivity, zero flags and accessibility by real quantifier
elimination. These are finite formulas: for example coverage is
$\forall\theta\ \forall x\in\Gamma_\theta\ \exists y$ in one of
the finitely many source charts with the stated projection. Each
positive neighbourhood and unit bound can likewise be represented by a
semialgebraic choice graph and its universal verification formula.

The existence proof in Appendix~\ref{app:foundations}, performed over
finite algebraic and rational-function extensions of the input field,
produces a certificate with such a finite encoding. An exhaustive search
therefore eventually finds a verified certificate. Invalid candidates are
rejected by the finite tests. This is an explicit input--output existence
algorithm; it is not a claim that the accompanying regression script
implements a practical resolution package. The direct maximum-function
algorithm above computes the spectral outputs without this enumeration.
Thus neither effectivity route depends on an unrepresented branch oracle,
and neither is inferred from a finite list of example calculations.
